\clearpage
\chapter*{Annexe}
\addcontentsline{toc}{chapter}{Annexe}
\markboth{Annexe}{Annexe}

\renewcommand{\thesection}{A.\arabic{section}}
\setcounter{section}{0}

\newcommand{\annexesection}[2]{%
    \refstepcounter{section}%
    \section*{\thesection\quad #1}%
    \label{#2}%
}

\annexesection{Contrat de sortie du moteur d'extraction structurelle}{annexe:contrat-moteur-extraction}

Le tableau~\ref{tab:annexe-rapport-extraction-sprint1} présente les principales sections du rapport d'extraction structurelle produit durant le Sprint 1.

\begingroup
\renewcommand{\arraystretch}{1.12}
\setlength{\tabcolsep}{4pt}
\begin{longtable}{|p{4.2cm}|p{8.8cm}|}
\caption{Contrat de sortie du moteur d'extraction structurelle}\label{tab:annexe-rapport-extraction-sprint1}\\
\hline
\textbf{Élément produit} & \textbf{Signification pour la migration} \\
\hline
Identité du projet & Permet de connaître la version Angular, le type d'environnement de travail et les outils utilisés. \\
\hline
Démarrage de l'application & Indique comment l'application est lancée et quelle stratégie d'initialisation est utilisée. \\
\hline
Routage & Sert à comprendre la navigation, les redirections, les routes protégées et les routes chargées à la demande. \\
\hline
Domaines fonctionnels & Aide à repérer les grandes parties métier candidates au futur découpage en micro-frontends. \\
\hline
Éléments partagés & Identifie les composants, services et éléments transverses à isoler ou mutualiser. \\
\hline
Dépendances internes & Met en évidence les couplages entre parties de l'application et les risques liés au découpage. \\
\hline
Intégration backend & Recense les appels aux API afin d'anticiper les dépendances entre front-end et services backend. \\
\hline
Sécurité & Distingue les zones publiques et protégées de l'application. \\
\hline
Gestion d'état & Identifie les mécanismes utilisés pour partager ou synchroniser les données dans l'application. \\
\hline
Complétude de l'analyse & Donne une indication sur la fiabilité et la couverture du rapport généré. \\
\hline
\end{longtable}
\endgroup

\annexesection{Contrat de sortie de l'agent de délimitation des domaines métier}{annexe:contrat-agent-domaines}

Le tableau~\ref{tab:annexe-sortie-agent-domaines} présente les principaux éléments produits par l'agent de délimitation des domaines métier durant le Sprint 2.

\begingroup
\renewcommand{\arraystretch}{1.12}
\setlength{\tabcolsep}{4pt}
\begin{longtable}{|p{4cm}|p{9.2cm}|}
\caption{Contrat de sortie de l'agent de délimitation des domaines métier}\label{tab:annexe-sortie-agent-domaines}\\
\hline
\textbf{Élément produit} & \textbf{Rôle} \\
\hline
Domaines métier candidats & Identifient les grandes parties fonctionnelles détectées dans le monolithe. \\
\hline
Candidats Remote & Indiquent les domaines pouvant être transformés progressivement en micro-frontends. \\
\hline
Éléments partagés & Regroupent les composants, services, modèles ou éléments transverses à mutualiser. \\
\hline
Risques de couplage & Signalent les dépendances pouvant rendre une séparation plus difficile. \\
\hline
Découpages différés & Indiquent les domaines dont la migration doit être reportée. \\
\hline
Niveau de confiance & Évalue la fiabilité globale des frontières proposées. \\
\hline
\end{longtable}
\endgroup

\annexesection{Contrat de sortie de l'agent de conception du plan directeur de migration}{annexe:contrat-agent-plan-directeur}

Le tableau~\ref{tab:annexe-sortie-agent-plan-directeur} présente les principaux éléments produits par l'agent de conception du plan directeur de migration durant le Sprint 2.

\begingroup
\renewcommand{\arraystretch}{1.12}
\setlength{\tabcolsep}{4pt}
\begin{longtable}{|p{4cm}|p{9.2cm}|}
\caption{Contrat de sortie de l'agent de conception du plan directeur de migration}\label{tab:annexe-sortie-agent-plan-directeur}\\
\hline
\textbf{Élément produit} & \textbf{Rôle} \\
\hline
Shell & Définit les responsabilités de l'application hôte, notamment le démarrage, le routage global et les fonctionnalités transverses. \\
\hline
Remotes & Décrit les domaines à migrer progressivement sous forme de micro-frontends. \\
\hline
Stratégie de routage & Précise les routes conservées par le Shell et les points d'entrée des Remotes. \\
\hline
Politique de partage & Identifie les éléments communs à mutualiser entre le Shell et les Remotes. \\
\hline
Ordre de migration & Propose une séquence de migration progressive et à faible risque. \\
\hline
Décisions architecturales & Justifient les choix retenus pour l'architecture cible. \\
\hline
Hypothèses et confiance & Indiquent les limites de la décision et le niveau de fiabilité du plan proposé. \\
\hline
\end{longtable}
\endgroup

\annexesection{Contrat de sortie de l'agent de génération de l'environnement de travail du projet et du Shell}{annexe:contrat-agent-workspace-shell}

Le tableau~\ref{tab:annexe-sortie-agent-workspace-shell} présente les principaux éléments produits par l'agent de génération de l'environnement de travail du projet et du Shell durant le Sprint 3.

\begingroup
\renewcommand{\arraystretch}{1.12}
\setlength{\tabcolsep}{4pt}
\begin{longtable}{|p{4cm}|p{9.2cm}|}
\caption{Contrat de sortie de l'agent de génération de l'environnement de travail du projet et du Shell}\label{tab:annexe-sortie-agent-workspace-shell}\\
\hline
\textbf{Élément produit} & \textbf{Rôle} \\
\hline
Plan du Shell & Décrit le nom du Shell, les routes contrôlées, les Remotes montés, les entrées distantes et les expositions partagées. \\
\hline
Plan de l'environnement de travail du projet & Définit le nom de l'environnement de travail, les applications, la librairie partagée et les commandes d'initialisation. \\
\hline
Configuration de bootstrap & Référence les fichiers de démarrage du Shell, notamment le point d'entrée principal, la configuration applicative et les routes. \\
\hline
Fichiers générés & Liste les fichiers produits pour le Shell, la configuration Angular, les environnements et les fichiers de fédération. \\
\hline
Commandes CLI & Présente les commandes permettant de recréer ou d'initialiser la structure Angular générée. \\
\hline
Chemin de l'environnement de travail du projet & Indique l'emplacement physique dans lequel les fichiers sont écrits. \\
\hline
Erreurs et avertissements & Signalent les problèmes rencontrés lors de la génération, de la normalisation ou de l'écriture des fichiers. \\
\hline
\end{longtable}
\endgroup

\annexesection{Contrat de sortie de l'agent de génération des squelettes Remote}{annexe:contrat-agent-remote-skeletons}

Le tableau~\ref{tab:annexe-sortie-agent-remote-skeletons} présente les principaux éléments produits par l'agent de génération des squelettes Remote durant le Sprint 3.

\begingroup
\renewcommand{\arraystretch}{1.12}
\setlength{\tabcolsep}{4pt}
\begin{longtable}{|p{4cm}|p{9.2cm}|}
\caption{Contrat de sortie de l'agent de génération des squelettes Remote}\label{tab:annexe-sortie-agent-remote-skeletons}\\
\hline
\textbf{Élément produit} & \textbf{Rôle} \\
\hline
Liste des Remotes & Décrit les applications Remote préparées, leurs noms de projet, leurs domaines et leurs routes de montage. \\
\hline
Fichiers générés & Liste les fichiers créés pour chaque Remote, notamment les fichiers de démarrage, de routes, de navigation, d'environnement et de fédération. \\
\hline
Projets intégrés & Indique les Remotes injectés dans la configuration globale de l'environnement de travail du projet. \\
\hline
Configuration des routes & Définit les routes locales et les points d'entrée exposés par chaque Remote. \\
\hline
Configuration de fédération & Prépare les fichiers nécessaires à l'exposition des routes Remote et à leur chargement par le Shell. \\
\hline
Chemin de l'environnement de travail du projet & Indique l'emplacement où les squelettes Remote sont écrits. \\
\hline
Erreurs et avertissements & Regroupe les erreurs de génération, d'intégration ou d'écriture détectées pendant l'exécution. \\
\hline
\end{longtable}
\endgroup

\renewcommand{\thesection}{\thechapter.\arabic{section}}